\section{在平面中嵌入图}\label{ch10:section-planar-graphs}

一个 (有限) 线性图 (linear graph) $G$ 是一个 Hausdorff 空间, 它可以表示成有限多个弧的并, 并且这些弧两两最多交于一个公共端点. 这些弧称为这个图的边 (edge), 弧的端点称为图的顶点 (vertex).

线性图常被数学家用来模拟现实生活中的现象, 然而我们只是把它们简单地看成一些有趣的空间, 在某种意义下作为简单闭曲线的推广.

对于任何一个图 (在不区别同胚的两个图的意义下), 只要列举出它的全部顶点, 并且指出哪些顶点的偶对有边将它们相连, 这个图便完全确定下来了.

\begin{example}\label{ch10:ex-complete-graph}
设 $G$ 包含 $n$ 个顶点, 并且 $G$ 中每对互异顶点都存在 $G$ 的一条边连接, 那么我们称 $G$ 是一个 $n$ 顶点完全图 (complete graph), 记为 $G_{n}$. \cref{ch10:fig-64-1} 中的几个图就满足上述条件. 注意, 前三个图都是 $\mathbb{R}^{2}$ 中的子空间, 第四个图是 $\mathbb{R}^{3}$ 中的子空间, 稍稍尝试之后可以发现这个图不能嵌入到 $\mathbb{R}^{2}$ 中, 稍后我们将证明这一点.
\end{example}

\par
\begin{figure}[htbp]
\centering
\includegraphics[width=0.7\textwidth]{ch10/fig-64-1.png}
\caption{完全图 $G_{2},G_{3},G_{4},G_{5}$.}\label{ch10:fig-64-1}
\end{figure}

\begin{example}\label{ch10:ex-utilities}
另外的例子是一个有趣的经典难题:``给定三所房子 $h_{1},h_{2},h_{3}$ 和三种功能 $g$ (通气), $w$ (通水), $e$ (通电), 问能否将每所房子与这三种功能都连接起来, 同时不让每两条连线相交?'' 将此问题归结为以下数学问题: 称为气水电图 (utilities graph) 的那个图 (\cref{ch10:fig-64-2}) 是否可以嵌入到 $\mathbb{R}^{2}$ 中? 稍稍尝试之后可以看出这是不可能的, 稍后会有证明.
\end{example}

\begin{definition}[theta space]\label{ch10:def-theta-space}
$\theta$ 空间 (theta space) $X$ 是指一个能够表示为三段弧 $A,B,C$ 的并的 Hausdorff 空间, 并且这三段弧中每两段都恰好相交于它们的两个端点. (这个空间 $X$ 显然同胚于希腊字母 $\theta$.)
\end{definition}

注意, 严格来说 $\theta$ 空间 $X$ 并不是线性图, 因为其中每两段弧的交多于一个公共端点. 但是我们能够把它表示为一个图, 只要把每一段弧 $A,B$ 和 $C$ 分拆成相交于一个公共端点的两段弧便可以了.

\begin{lemma}\label{ch10:lemma-theta-space}
设 $X$ 是一个 $\theta$ 空间, 并且它还是 $S^{2}$ 的子空间, $A,B$ 和 $C$ 是弧, 它们的并恰好是空间 $X$. 那么 $X$ 将 $S^{2}$ 分成三个分支, 它们的边界分别是 $A\cup B$, $B\cup C$ 和 $A\cup C$. 以 $A\cup B$ 作为边界的分支就是 $S^{2}-A\cup B$ 的分支之一.
\end{lemma}

\begin{proof}
设 $a$ 和 $b$ 是弧 $A,B$ 和 $C$ 的端点. 考虑简单闭曲线 $A\cup B$, 它将 $S^{2}$ 分成两个分支 $U$ 和 $U'$, 并且这两个分支都是 $S^{2}$ 中的开集, 具有公共边界 $A\cup B$. 参见 \cref{ch10:fig-64-3}.
\end{proof}

\par
\begin{figure}[htbp]
\centering
\includegraphics[width=0.8\textwidth]{ch10/fig-64-2-3.png}
\caption{左: 气水电图; 右: $\theta$ 空间的三个分支.}\label{ch10:fig-64-2-3}
\end{figure}

\begin{proof}[continued]
因为空间 $C-a-b$ 是连通的, 所以包含在其中的一个分支中, 不妨设为 $U'$. 现在考虑两个空间 $\overline{U}=U\cup A\cup B$ 和 $C$, 它们都是连通的. 因为 $C$ 是一段弧, $\overline{U}$ 的补是连通集 $U'$, 所以两者都不能分割 $S^{2}$. 又因为这两个集合的交只有两个点 $a$ 和 $b$, 于是根据 \cref{ch10:thm-two-arcs-nonseparating}, 它们的并将 $S^{2}$ 分割成两个分支 $V$ 和 $W$. 由此可见 $S^{2}-(A\cup B\cup C)$ 可以表示成三个无交连通集 $U,V$ 和 $W$ 的并. 因为这三个集合都是 $S^{2}$ 中的开集, 所以它们就是 $S^{2}-(A\cup B\cup C)$ 的分支, 其中 $U$ 的边界就是 $A\cup B$. 类似地, 其他两个分支的边界分别为 $B\cup C$ 和 $A\cup C$.
\end{proof}

\begin{theorem}\label{ch10:thm-utilities-nonplanar}
设 $X$ 是气水电图. 则 $X$ 不能嵌入平面.
\end{theorem}

\begin{proof}
如果 $X$ 能嵌入到平面中, 那么就能嵌入到 $S^{2}$ 中, 因此我们假设 $X$ 是 $S^{2}$ 的一个子空间, 由此推出矛盾.

我们使用例 2 中的记号, $g,w,e,h_{1},h_{2}$ 和 $h_{3}$ 是 $X$ 的顶点. 设 $A,B$ 和 $C$ 是 $X$ 中的弧:
\[
A=gh_{1}w,\qquad B=gh_{2}w,\qquad C=gh_{3}w.
\]
这三段弧中的任意两段只交于端点 $g$ 和 $w$, 因此 $Y=A\cup B\cup C$ 是一个 $\theta$ 空间. 空间 $Y$ 将 $S^{2}$ 分割为三个分支 $U,V$ 和 $W$, 它们的边界分别为 $A\cup B$, $B\cup C$ 和 $A\cup C$. 参见 \cref{ch10:fig-64-4}.

$X$ 的顶点 $e$ 属于这三个分支之一, 因此 $X$ 中的弧 $eh_{1},eh_{2}$ 和 $eh_{3}$ 都包含在这个分支的闭包中. 这个分支不可能是 $U$, 因为 $\overline{U}$ 包含在 $U\cup A\cup B$ 中, 这个集合不包含点 $h_{3}$. 类似地, 包含 $e$ 的分支也不可能是 $V$ 和 $W$, 因为 $\overline{V}$ 不包含 $h_{1}$, $\overline{W}$ 不包含 $h_{2}$. 从而推出矛盾.
\end{proof}

\par
\begin{figure}[htbp]
\centering
\includegraphics[width=0.5\textwidth]{ch10/fig-64-4.png}
\caption{气水电图的 $\theta$ 子图 $Y=A\cup B\cup C$.}\label{ch10:fig-64-4}
\end{figure}

\begin{lemma}\label{ch10:lemma-k4-divides}
设 $X$ 是 $S^{2}$ 中以 $a_{1},a_{2},a_{3}$ 和 $a_{4}$ 为顶点的完全图. 则 $X$ 将 $S^{2}$ 分成四个分支. 设这四个分支的边界分别为 $X_{1},X_{2},X_{3}$ 和 $X_{4}$, 则每一个 $X_{i}$ 就是 $X$ 中不以 $a_{i}$ 为顶点的那些边的并.
\end{lemma}

\begin{proof}
设 $Y$ 是 $X$ 中异于弧 $a_{2}a_{4}$ 的所有弧的并. 我们可以按以下方式将 $Y$ 表示成一个 $\theta$ 空间:
\[
A=a_{1}a_{2}a_{3},\qquad B=a_{1}a_{3},\qquad C=a_{1}a_{4}a_{3}.
\]
参见 \cref{ch10:fig-64-5}. 弧 $A,B$ 和 $C$ 在端点 $a_{1}$ 和 $a_{3}$ 处相交, 并且它们的并就是 $Y$.
\end{proof}

\par
\begin{figure}[htbp]
\centering
\includegraphics[width=0.5\textwidth]{ch10/fig-64-5.png}
\caption{$G_{4}$ 中的 $\theta$ 子图 $Y$.}\label{ch10:fig-64-5}
\end{figure}

\begin{proof}[continued]
空间 $Y$ 将 $S^{2}$ 分割成三个分支 $U,V$ 和 $W$, 它们的边界分别是 $A\cup B$, $B\cup C$ 和 $A\cup C$. 空间 $a_{2}a_{4}-a_{2}-a_{4}$ 是连通的, 必包含于这三个分支之一. 但它不包含于 $U$ 中, 因为 $A\cup B$ 不包含 $a_{4}$. 同时它也不包含于 $V$, 因为 $B\cup C$ 不包含 $a_{2}$. 从而它只能包含于 $W$.

$\overline{U}\cup\overline{V}$ 是连通的, 因为 $\overline{U}$ 和 $\overline{V}$ 是连通的, 并且有非空的交 $B$. 进而, 集合 $\overline{U}\cup\overline{V}$ 不分割 $S^{2}$, 因为它的补是 $W$. 类似地, 弧 $a_{2}a_{4}$ 是连通的, 不分割 $S^{2}$. 集合 $a_{2}a_{4}$ 与 $\overline{U}\cup\overline{V}$ 仅交于 $a_{2}$ 和 $a_{4}$. 根据 \cref{ch10:thm-two-arcs-nonseparating}, $a_{2}a_{4}\cup\overline{U}\cup\overline{V}$ 将 $S^{2}$ 分割为两个分支 $W_{1}$ 和 $W_{2}$, 所以 $S^{2}-Y$ 是四个无交连通集合 $U,V,W_{1}$ 和 $W_{2}$ 的并. 因为这几个集合都是开集, 所以它们便是 $S^{2}-Y$ 的诸分支.

在这些分支中, 不妨将其中之一设为 $U$, 并以图 $A\cup B=X_{4}$ 作为边界. 通过类似的讨论可见, 其他三个分支分别将 $X_{1},X_{2}$ 和 $X_{3}$ 作为边界.
\end{proof}

\begin{theorem}\label{ch10:thm-k5-nonplanar}
五个顶点的完全图不能嵌入平面.
\end{theorem}

\begin{proof}
设 $G$ 是 $S^{2}$ 的子空间, 并且是以 $a_{1},a_{2},a_{3},a_{4}$ 和 $a_{5}$ 为五个顶点的完全图. 令 $X$ 表示 $G$ 的诸边中不以 $a_{5}$ 为顶点的那些边之并, 从而 $X$ 是四个顶点的完全图. 于是 $X$ 将 $S^{2}$ 分割成四个分支, 这四个分支的边界分别为 $X_{1},X_{2},X_{3}$ 和 $X_{4}$, 其中 $X_{i}$ 是不以 $a_{i}$ 为顶点的那些边构成的. 点 $a_{5}$ 必属于这四个分支之一, 于是那些以 $a_{5}$ 为顶点的边之并构成的连通集合
\[
a_{1}a_{5}\cup a_{2}a_{5}\cup a_{3}a_{5}\cup a_{4}a_{5}
\]
也必包含在这个分支的闭包中. 因此 $a_{1},\cdots,a_{4}$ 都属于这个分支的边界. 但这是不可能的, 因为没有哪个 $X_{i}$ 能够包含所有四个顶点 $a_{1},\cdots,a_{4}$, 于是产生了矛盾.
\end{proof}

根据这些定理可以得到: 若图 $G$ 包含一个子图是气水电图或者是五个顶点的完全图, 则 $G$ 不能嵌入平面. 逆命题也是正确的, 它是一个著名的定理, 由 Kuratowski 证得, 但证明并不容易.

\begin{exercise}\label{ch10:ex-planar-graph-1}
设 $X$ 是一个空间, 可表示为有限多段弧 $A_{1},\cdots,A_{n}$ 的并, 这些弧中的每一对最多交于一个公共端点.
\end{exercise}